\section{What the children make available to a categorical root}
\label{sec:categorical-root}

Fix a binary RHM grammar of depth four and a TTN whose tree agrees with the
grammar. The root label is uniform on the $n$ classes. Each grammar symbol
chooses one of its $m$ ordered productions uniformly, independently of the
choices below it. Write $H_L,H_R\in\mathbb R^r$ for the representations
of the two eight-token subtrees and $Z=H_L\otimes H_R\in\mathbb R^{r^2}$
for their ordered product. In the repository's \texttt{TreeTN}, the root
maps this product directly to logits $CZ$, with $C\in\mathbb R^{n\times
r^2}$ and no bias. All formulas below apply to any fixed choice of the
lower cores, including a snapshot of jointly trained weights. For time
evolution with $Z$ fixed, the optimizer specified below is population
gradient flow. The repository trains with minibatch Adam.

\subsection{The population gradient of the actual classifier}

Let $e_Y$ denote the one-hot root label, let $p=\mathbf{1}/n$, and set
$D_j=\mathbb E[Z\mid Y=j]$ and $\bar D=\mathbb E Z$.
For cross-entropy, differentiating the log partition function gives
\begin{align}
 {\cal L}_{\rm CE}(C)
 &=\mathbb E\biggl[\log\sum_j e^{(CZ)_j} -(CZ)_Y\biggr],\\
 -\nabla_C{\cal L}_{\rm CE}(C)
 &=\mathbb E\big[(e_Y-\operatorname{softmax}(CZ))Z^\top\big].
 \label{eq:cat-full-gradient}
\end{align}
Resetting the root to zero makes its prediction $p$ for every tree. Its
initial descent direction $B$ therefore has rows
\begin{equation}
 B=\mathbb E[(e_Y-p)Z^\top],\qquad
 B_{j,:}=\frac{1}{n}(D_j-\bar D)^\top .
 \label{eq:cat-zero-gradient}
\end{equation}
If $B=0$, convexity of cross-entropy in the root weights makes $C=0$ a
global minimizer with loss $\log n$ on these frozen children. If $B\ne0$,
moving the root to $\eta B$ improves the loss by
$\eta\|B\|_F^2+O(\eta^2)$ for small $\eta>0$. Thus $B$ decides
\emph{exactly} whether a bilinear root can improve population
cross-entropy over the uniform predictor. This test compares full
matrices. One vanishing row or one vanishing root contrast establishes
only zero initial drive along that direction.

For a centered class contrast $c\in\mathbb R^n$ with $\sum_j c_j=0$,
write $R_j(u,v)=\Pr(U=u,V=v\mid Y=j)$ for the root rule probabilities
and $K_c(u,v)=n^{-1}\sum_j c_jR_j(u,v)$ for the signed rule table.
Let $M_L,M_R$ have columns $\mathbb E[H_L\mid U=u]$ and
$\mathbb E[H_R\mid V=v]$, respectively. Given the child symbols
$(u,v)$, the two subtrees expand independently. Averaging first over the
descendants and then over the root production gives
\begin{equation}
 \mathbb E[c_YH_LH_R^\top]
 =\sum_{u,v}K_c(u,v)\,M_L(:,u)M_R(:,v)^\top
 =M_LK_c M_R^\top .
 \label{eq:cat-contrast}
\end{equation}
Flattening this matrix gives $c^\top B$ in the order used to define $Z$.
For a singular term $\sigma a b^\top$ of $K_c$, its contribution is
$\sigma(M_La)(M_Rb)^\top$. A child that annihilates its side of this term
removes it from the drive. Other terms can survive or cancel in the sum.
When $K_c$ has zero row and column sums, collapsing either child's
conditional means to one common vector erases the entire contrast from
the initial drive. A generic grammar need not have these vanishing
marginals. For instance, a legal root table can encode the root class in
the left symbol alone, in which case even a constant nonzero right bond
can suffice.

\subsection{Exact moments without enumerating the leaf strings}

At each node $i$ with grammar symbol $s$ define
$\mu_i(s)=\mathbb E[H_i\mid s]$ and
$S_i(s)=\mathbb E[H_i H_i^\top\mid s]$. Use the tensor convention
$(x\otimes y)_{ab}=x_a y_b$. At leaf $i$, whose embedding of $s$ is
$h_i(s)$, the initial values are $h_i(s)$ and $h_i(s)h_i(s)^\top$.
At an internal node write the flattened core as $W_i$, so
$H_i=W_i(H_{iL}\otimes H_{iR})$. If the $k$th legal production of $s$
is $(u_k,v_k)$, conditional independence gives
\begin{align}
 d_i(s)&=\frac1m\sum_{k=1}^m \mu_{iL}(u_k)\otimes\mu_{iR}(v_k),\\
 F_i(s)&=\frac1m\sum_{k=1}^m S_{iL}(u_k)\otimes S_{iR}(v_k),\\
 \mu_i(s)&=W_i d_i(s),\qquad S_i(s)=W_i F_i(s)W_i^\top .
 \label{eq:cat-moments}
\end{align}
The second line uses $(x\otimes y)(x\otimes y)^\top
=(xx^\top)\otimes(yy^\top)$ before taking the conditional expectation.
At the root, stop before multiplying by its core. Then $D_j=d_{\rm
root}(j)$ and $F=n^{-1}\sum_j F_{\rm root}(j)=\mathbb E[ZZ^\top]$.
These contractions average only legal productions, with their true
probabilities. They apply to level-tied and untied weights because the
conditioning fixes the weights. A TTN that permutes leaves across
grammar subtrees requires a different contraction.

For any matrix perturbation $E$ of the root, the Hessian at zero obeys
\begin{equation}
 \nabla^2 {\cal L}_{\rm CE}(0)[E]
 =(\operatorname{diag}p-pp^\top)E F .
\end{equation}
This follows by differentiating the softmax in
\eqref{eq:cat-full-gradient}. At a nonzero root the softmax Jacobian
depends on the particular tree and stays inside the expectation. Thus
the two moments determine the exact initial gradient and curvature.
Evaluating the full nonlinear gradient at a trained root generally
requires more of the distribution of $Z$.

\subsection{Fitting a root above frozen children}

To quantify what a frozen representation can predict at arbitrary root
weights, take a scalar target $y=c_Y$ and fit the same bias-free root
with population squared error. With $b=\mathbb E[Zy]$, expansion gives
\begin{equation}
 2{\cal L}_2(w)=\mathbb E[y^2]-2w^\top b+w^\top Fw .
\end{equation}
For $v\in\ker F$, $\mathbb E[(v^\top Z)^2]=0$, so $v^\top b=0$ as well.
Consequently $b$ lies in the range of $F$, $w_*=F^\dagger b$ is a
minimizer, and completing the square gives
\begin{equation}
 2{\cal L}_2(w)=\mathbb E[y^2]-b^\top F^\dagger b
                  +(w-w_*)^\top F(w-w_*) .
 \label{eq:cat-oracle}
\end{equation}
Here $F^\dagger$ is the Moore--Penrose pseudoinverse. For a nonzero contrast,
the fraction of its variance explained by the root is
$b^\top F^\dagger b/\mathbb E[y^2]$. It is invariant under any
invertible change of basis on either child bond. Indeed, such a change preserves the set of
scalar predictions $w^\top Z$ after reparametrizing $w$. The Euclidean
gradient and its speed depend on that parametrization.

For a fresh root $w(0)=0$, gradient flow obeys $\dot w=b-Fw$.
Diagonalize $F$ with orthonormal eigenvectors $v_k$ and positive
eigenvalues $\lambda_k$. Each component obeys
\begin{equation}
 v_k^\top w(t)=\frac{v_k^\top b}{\lambda_k}
       (1-e^{-\lambda_k t}),\qquad
 {\cal L}_2(w(t))-{\cal L}_2(w_*)
 =\frac12\sum_{\lambda_k>0}
 \frac{(v_k^\top b)^2}{\lambda_k}e^{-2\lambda_k t}.
\end{equation}
A nonzero component reaches half its final amplitude at
$\log 2/\lambda_k$. This separates a better attainable fit from a
shorter learning time in fixed coordinates. These are regression
half-times for freshly trained roots at different lower-core snapshots.
They do not give hitting times for the jointly trained categorical model.

\paragraph{Why a shared factor does not imply a lag.}
In the symmetric two-mode model $f_i=a b_i z_i$ with equal targets and
equal private initial conditions, exchanging indices $1,2$ leaves the
loss and initialization invariant. Uniqueness of gradient flow then
implies $b_1(t)=b_2(t)$ and $z_1(t)=z_2(t)$, so the two modes hit every
common threshold together despite sharing the trainable factor $a$.
More generally a change $H_L\mapsto \alpha H_L$, $C\mapsto C/\alpha$
preserves every categorical prediction while changing the apparent
size of the child feature. A quantitative time-gap theorem must specify
a distribution, a target contrast, the initialization, and the optimizer's
metric, as well as the intervention that tests necessity.
